\section{Introducción}

En este tercer hito se aborda la implementación integral del módulo de usuario, tomando como punto de partida el modelo de base de datos definido en el hito anterior. A partir del diagrama Entidad–Relación se proyectan los tres módulos principales del sistema, usuarios, productos y préstamos, siendo el módulo de usuario el primero en implementarse y, al mismo tiempo, el más relevante, ya que establece la base tecnológica y conceptual sobre la cual se construirán los siguientes hitos.

Este módulo contempla las tablas \texttt{users} y \texttt{roles}, encargadas de gestionar la información básica de quienes interactúan con el laboratorio y de definir los diferentes niveles de permisos y responsabilidades dentro del sistema. Su correcta implementación resulta crítica, pues determina el flujo de acceso y autenticación de los usuarios, y a la vez sienta un patrón de diseño para los futuros servicios asociados a los modelos de productos y préstamos. De este modo, se garantiza coherencia en la arquitectura modular, consistencia en la forma de exponer los servicios vía API y una base sólida para incorporar buenas prácticas de desarrollo y documentación a lo largo de todo el proyecto.
